% ==============================================================================
% Paper 84 (DE): Die Fontaine-Mazur-Vermutung:
%                p-adische Galois-Darstellungen und Modularität
% Autor: Michael Fuhrmann
% Projekt: Specialist Mathematics Project, Build 167
% Datum: 2026-03-12
% ==============================================================================

\documentclass[12pt,a4paper]{amsart}

\usepackage{amsmath,amssymb,amsthm,mathtools,enumitem,booktabs,array,hyperref}
\usepackage[utf8]{inputenc}
\usepackage[T1]{fontenc}
\usepackage[ngerman]{babel}
\usepackage{geometry}
\geometry{margin=2.5cm}

% --------------------------------------------------------------------------
% Theorem-Umgebungen (Deutsch)
% --------------------------------------------------------------------------
\newtheorem{theorem}{Theorem}[section]
\newtheorem{lemma}[theorem]{Lemma}
\newtheorem{corollary}[theorem]{Korollar}
\newtheorem{proposition}[theorem]{Proposition}
\newtheorem{satz}[theorem]{Satz}
\newtheorem{korollar}[theorem]{Korollar}
\theoremstyle{definition}
\newtheorem{definition}[theorem]{Definition}
\newtheorem{definition_de}[theorem]{Definition}
\newtheorem{example}[theorem]{Beispiel}
\theoremstyle{remark}
\newtheorem{remark}[theorem]{Bemerkung}
\newtheorem{bemerkung}[theorem]{Bemerkung}
\newtheorem{conjecture}[theorem]{Vermutung}
\newtheorem{vermutung}[theorem]{Vermutung}

% --------------------------------------------------------------------------
% Makros
% --------------------------------------------------------------------------
\newcommand{\Gal}{\mathrm{Gal}}
\newcommand{\GL}{\mathrm{GL}}
\newcommand{\SL}{\mathrm{SL}}
\newcommand{\Qp}{\mathbb{Q}_p}
\newcommand{\Zp}{\mathbb{Z}_p}
\newcommand{\Qbar}{\overline{\mathbb{Q}}}
\newcommand{\GQ}{G_{\mathbb{Q}}}
\newcommand{\BdR}{B_{\mathrm{dR}}}
\newcommand{\Bcris}{B_{\mathrm{cris}}}
\newcommand{\Bst}{B_{\mathrm{st}}}
\newcommand{\DdR}{D_{\mathrm{dR}}}
\newcommand{\Dcris}{D_{\mathrm{cris}}}
\newcommand{\Dst}{D_{\mathrm{st}}}
\newcommand{\Frob}{\mathrm{Frob}}
\newcommand{\rec}{\mathrm{rec}}
\newcommand{\Z}{\mathbb{Z}}
\newcommand{\Q}{\mathbb{Q}}
\newcommand{\R}{\mathbb{R}}
\newcommand{\C}{\mathbb{C}}
\newcommand{\F}{\mathbb{F}}
\newcommand{\N}{\mathbb{N}}

\hypersetup{
  colorlinks=true,
  linkcolor=blue,
  citecolor=blue,
  urlcolor=blue,
  pdftitle={Die Fontaine-Mazur-Vermutung},
  pdfauthor={Michael Fuhrmann}
}

% ==============================================================================
\begin{document}
% ==============================================================================

\title{Die Fontaine-Mazur-Vermutung:\\
  $p$-adische Galois-Darstellungen und Modularität}

\author{Michael Fuhrmann}
\address{Specialist Mathematics Project, Build 167}
\email{specialist-maths@project.local}
\date{2026-03-12}

\begin{abstract}
Die Fontaine-Mazur-Vermutung, formuliert 1995 von Fontaine und Mazur \cite{FM1995},
besagt, dass jede irreduzible geometrische $p$-adische Galois-Darstellung
$\rho\colon \GQ \to \GL_n(\Qp)$
aus der Geometrie stammt und -- im Fall $n=2$ -- modular ist, d.\,h.\ einer klassischen
Modulform entspricht.
Wir entwickeln die notwendigen Grundlagen der $p$-adischen Hodge-Theorie (Periodenringe
$\BdR$, $\Bcris$, $\Bst$), erläutern die Hierarchie der zulässigen Darstellungen
(de Rham, kristallin, semistabil, Hodge-Tate) und formulieren die Vermutung präzise.
Anschließend besprechen wir die wichtigsten bewiesenen Fälle:
Wiles' Satz über elliptische Kurven ($n=2$, Gewicht $2$),
Kisins Arbeit (2009) zu potenziell semistabilen Darstellungen,
Emertons Methode der komplettierten Kohomologie und
Taylors potenzieller Modularitätssatz.
Für $n \geq 3$ bleibt die Vermutung vollständig offen.
\end{abstract}

\maketitle

\tableofcontents

% ==============================================================================
\section{Einleitung}
% ==============================================================================

Ein zentrales Problem der modernen Zahlentheorie ist es, alle
$\ell$-adischen Galois-Darstellungen
\[
  \rho \colon \GQ \;:=\; \Gal(\Qbar/\Q) \;\longrightarrow\; \GL_n(\overline{\Q}_\ell)
\]
zu klassifizieren: Welche kommen \emph{aus der Geometrie} (d.\,h.\ treten als
Unterquotienten der étalen Kohomologie einer glatten projektiven Varietät über $\Q$ auf)?
Welche sind \emph{automorph} (d.\,h.\ korrespondieren zu automorphen Darstellungen von
$\GL_n$ über dem Adele-Ring $\mathbb{A}_\Q$)?

Die Fontaine-Mazur-Vermutung \cite{FM1995} gibt eine präzise Antwort auf diese Frage für
$p$-adische Galois-Darstellungen: Eine irreduzible, geometrische, ungerade ($n=2$)
$p$-adische Darstellung soll modular sein.

\subsection{Historischer Hintergrund}

Die systematische Untersuchung $p$-adischer Galois-Darstellungen wurde durch Fontaines
Entwicklung der $p$-adischen Hodge-Theorie in den 1980er und 1990er Jahren ermöglicht
\cite{Fontaine1979,Fontaine1982,Fontaine1994}.
Fontaine führte eine Hierarchie ``schöner'' Darstellungen ein:
Hodge-Tate $\supset$ de Rham $\supset$ semistabil $\supset$ kristallin,
und ordnete jeder Klasse einen linearen algebraischen Gegenstand zu
(gefiltertes $\phi$-Modul, $(\phi,N)$-Modul usw.).

1995 formulierten Fontaine und Mazur \cite{FM1995} ihre weitreichende Vermutung,
die besagt, dass jede irreduzible \emph{geometrische} Darstellung
(d.\,h.\ de Rham bei $p$ und fast überall unverzweigt) aus der Geometrie stammt
und damit automorph ist.

Der $n=2$-Fall wurde in seinen wichtigsten Instanzen durch Wiles \cite{Wiles1995} und
Taylor-Wiles \cite{TW1995} bewiesen (halbstabile elliptische Kurven; vollständig für
alle elliptischen Kurven durch Breuil-Conrad-Diamond-Taylor \cite{BCDT2001}),
und in viel größerer Allgemeinheit durch Kisin \cite{Kisin2009} und Emerton
\cite{Emerton2010}.

\subsection{Aufbau des Aufsatzes}

Abschnitt~\ref{sec:periodenringe} behandelt die Periodenringe der $p$-adischen
Hodge-Theorie.
Abschnitt~\ref{sec:darstellungen} erläutert zulässige Darstellungen.
Abschnitt~\ref{sec:galoisdarstellungen} beschreibt die globale Theorie.
Abschnitt~\ref{sec:vermutung} formuliert die Fontaine-Mazur-Vermutung.
Abschnitt~\ref{sec:langlands} verbindet die Vermutung mit dem Langlands-Programm.
Abschnitt~\ref{sec:wiles} bespricht Wiles' Satz.
Abschnitt~\ref{sec:kisin} behandelt Kisins Arbeit.
Abschnitt~\ref{sec:emerton} beschreibt Emertons komplettierende Kohomologie.
Abschnitt~\ref{sec:taylor} erläutert Taylors potenzielle Modularität.
Abschnitt~\ref{sec:offene_probleme} listet offene Probleme auf.
Abschnitt~\ref{sec:fazit} schließt den Aufsatz ab.

% ==============================================================================
\section{$p$-adische Hodge-Theorie: Periodenringe}\label{sec:periodenringe}
% ==============================================================================

Sei $p$ eine Primzahl und $K/\Qp$ eine endliche Erweiterung mit absoluter
Galois-Gruppe $G_K = \Gal(\overline{K}/K)$.

\subsection{Der de-Rham-Periodenring $\BdR$}

\begin{definition_de}[Fontaine \cite{Fontaine1982}]
  Sei $\widetilde{\mathbf{E}}^+ = \varprojlim_{x \mapsto x^p} \mathcal{O}_{\overline{K}}/p$
  der \emph{Tilt} von $\mathcal{O}_{\overline{K}}$. Der Ring
  $A_{\inf}(\mathcal{O}_{\overline{K}}) = W(\widetilde{\mathbf{E}}^+)$
  ist der Witt-Ring von $\widetilde{\mathbf{E}}^+$. Es gibt einen surjektiven
  Ringhomomorphismus $\theta\colon A_{\inf} \to \mathcal{O}_{\mathbb{C}_p}$,
  dessen Kern ein Hauptideal ist, erzeugt von einem Nicht-Nullteiler $\xi$.

  Der \emph{de-Rham-Periodenring} ist
  \[
    \BdR \;:=\; \varprojlim_{n} A_{\inf}[\tfrac{1}{p}] / (\ker\theta)^n,
  \]
  ein vollständiger diskreter Bewertungskörper mit Restklassenkörper $\mathbb{C}_p$
  und Uniformisierender $t = \log([\varepsilon])$, wo
  $\varepsilon = (1, \zeta_p, \zeta_{p^2},\ldots)$ ein kompatibles System von
  $p$-Potenzen von Einheitswurzeln ist.
\end{definition_de}

Der Ring $\BdR$ trägt eine Filtration
$\mathrm{Fil}^i \BdR = t^i \BdR^+$
und eine $G_K$-Wirkung; die Invarianten erfüllen $\BdR^{G_K} = K$.

\subsection{Der kristalline Periodenring $\Bcris$}

\begin{definition_de}
  Der \emph{kristalline Periodenring} $\Bcris$ ist die $p$-adische Vervollständigung von
  \[
    A_{\mathrm{cris}} = \left\{ \sum_{n=0}^\infty a_n \frac{\xi^n}{n!} \,:\,
    a_n \in A_{\inf}, \; a_n \to 0 \right\},
  \]
  invertiert bei $p$. Er trägt einen Frobenius $\phi$, eine Filtration (induziert von
  $\BdR$) und eine $G_K$-Wirkung. Die Invarianten erfüllen
  $\Bcris^{G_K} = K_0$ (der maximale unverzweigte Unterkörper von $K$).
\end{definition_de}

\subsection{Der semistabile Periodenring $\Bst$}

\begin{definition_de}
  Sei $\pi \in K$ ein Uniformisierer und $\ell = \log([\widetilde{\pi}]) \in \Bcris$.
  Der \emph{semistabile Periodenring} ist
  \[
    \Bst \;:=\; \Bcris[\ell].
  \]
  Er trägt den Frobenius $\phi$ (Fortsetzung von $\Bcris$), den
  \emph{Monodromie-Operator} $N\colon \Bst \to \Bst$ mit $N(\ell) = -1$,
  $N|_{\Bcris} = 0$ sowie eine $G_K$-Wirkung.
  Es gilt $\Bcris \subset \Bst \subset \BdR$.
\end{definition_de}

\begin{bemerkung}
  Das zyklotemische Zeichen $\chi_{\mathrm{zyk}}\colon G_K \to \Zp^\times$ erfüllt
  $g \cdot t = \chi_{\mathrm{zyk}}(g) \cdot t$ für das Element $t \in \BdR^+$.
\end{bemerkung}

% ==============================================================================
\section{Zulässige Darstellungen}\label{sec:darstellungen}
% ==============================================================================

Sei $V$ ein endlichdimensionaler $\Qp$-Vektorraum mit stetiger $G_K$-Wirkung
(eine $p$-adische Darstellung von $G_K$).

\begin{definition_de}
  \begin{enumerate}[label=(\roman*)]
    \item $V$ heißt \emph{Hodge-Tate}, falls
      $V \otimes_{\Qp} \mathbb{C}_p \cong \bigoplus_i \mathbb{C}_p(i)^{h_i}$
      als $G_K$-Moduln (mit $h_i \geq 0$).
    \item $V$ heißt \emph{de Rham}, falls
      $\DdR(V) := (V \otimes_{\Qp} \BdR)^{G_K}$
      die $K$-Dimension $\dim_{\Qp} V$ hat.
    \item $V$ heißt \emph{semistabil}, falls
      $\Dst(V) := (V \otimes_{\Qp} \Bst)^{G_K}$
      die $K_0$-Dimension $\dim_{\Qp} V$ hat.
    \item $V$ heißt \emph{kristallin}, falls
      $\Dcris(V) := (V \otimes_{\Qp} \Bcris)^{G_K}$
      die $K_0$-Dimension $\dim_{\Qp} V$ hat.
  \end{enumerate}
\end{definition_de}

\begin{satz}[Fontaine {\cite[Theorem A]{Fontaine1994}}]
  Es gibt eine strikte Inklusion voller Unterkategorien von $p$-adischen Darstellungen
  von $G_K$:
  \[
    \{ \text{kristallin} \} \;\subsetneq\; \{ \text{semistabil} \} \;\subsetneq\;
    \{ \text{de Rham} \} \;\subsetneq\; \{ \text{Hodge-Tate} \}.
  \]
\end{satz}

\begin{satz}[$C_{\mathrm{dR}}$-Vermutung $=$ Theorem, Berger \cite{Berger2002}]
  Jede de-Rham-Darstellung von $G_K$ ist potenziell semistabil, d.\,h.\ sie wird über
  einer endlichen Erweiterung von $K$ semistabil.
\end{satz}

\begin{bemerkung}
  Bergers Beweis \cite{Berger2002} nutzt die Theorie der $p$-adischen
  Differentialgleichungen, insbesondere das Monodromieproblem von Fontaine und die
  Theorie der Wach-Moduln.
\end{bemerkung}

\begin{definition_de}
  Die \emph{Hodge-Tate-Gewichte} einer de-Rham-Darstellung $V$ der Dimension $n$ sind
  die ganzen Zahlen $h_1 \leq h_2 \leq \cdots \leq h_n$, für die
  $\mathrm{gr}^{h_j} \DdR(V) \neq 0$.
\end{definition_de}

\begin{example}
  Der Tate-Modul $T_p E$ einer elliptischen Kurve $E/\Q_p$ mit guter Reduktion ist
  kristallin mit Hodge-Tate-Gewichten $\{0, 1\}$.
  Das zyklotomische Zeichen $\chi_{\mathrm{zyk}}$ ist kristallin mit
  Hodge-Tate-Gewicht $1$.
\end{example}

% ==============================================================================
\section{$p$-adische Galois-Darstellungen: Globale Theorie}\label{sec:galoisdarstellungen}
% ==============================================================================

Wir arbeiten nun global mit $\GQ = \Gal(\Qbar/\Q)$.

\begin{definition_de}
  Eine stetige Darstellung $\rho\colon \GQ \to \GL_n(\Qp)$ heißt:
  \begin{enumerate}[label=(\roman*)]
    \item \emph{unverzweigt bei $\ell \neq p$}, falls $\rho|_{I_\ell}$ trivial ist;
    \item \emph{geometrisch} (im Sinne von Fontaine-Mazur), falls sie de Rham bei $p$
      und fast überall unverzweigt ist;
    \item \emph{von geometrischem Ursprung}, falls sie als Unterquotient von
      $H^i_{\text{\'et}}(X_{\Qbar}, \Qp)(j)$ für eine glatte projektive Varietät $X/\Q$,
      ein $i$ und eine Tate-Drehung $j$ auftritt.
  \end{enumerate}
\end{definition_de}

\begin{example}[Tate-Modul einer elliptischen Kurve]
  Für eine elliptische Kurve $E/\Q$ ist der $p$-adische Tate-Modul
  $V_p E = T_p E \otimes_{\Zp} \Qp$
  eine irreduzible $2$-dimensionale $p$-adische Galois-Darstellung.
  Sie ist geometrisch (de Rham bei $p$, fast überall unverzweigt) und von
  geometrischem Ursprung (aus $H^1_{\text{\'et}}(E_{\Qbar}, \Qp)$).
\end{example}

\begin{example}[Galois-Darstellung einer Modulform]
  Sei $f = \sum_{n=1}^\infty a_n q^n \in S_k(\Gamma_0(N), \chi)$ eine normalisierte
  Neuform vom Gewicht $k \geq 2$. Deligne \cite{Deligne1971} konstruierte eine
  $2$-dimensionale $\lambda$-adische Darstellung
  \[
    \rho_f \colon \GQ \to \GL_2(\mathcal{O}_{f,\lambda}),
  \]
  die de Rham bei $p$ mit Hodge-Tate-Gewichten $\{0, k-1\}$ ist.
\end{example}

% ==============================================================================
\section{Die Fontaine-Mazur-Vermutung}\label{sec:vermutung}
% ==============================================================================

\begin{vermutung}[Fontaine-Mazur \cite{FM1995}]\label{verm:FM}
  Sei $\rho\colon \GQ \to \GL_n(\Qp)$ eine irreduzible, geometrische $p$-adische
  Galois-Darstellung. Dann gilt:
  \begin{enumerate}[label=(\alph*)]
    \item \textbf{(Motivisch)} $\rho$ ist von geometrischem Ursprung: Sie tritt in der
      $p$-adischen étalen Kohomologie einer glatten projektiven Varietät über $\Q$ auf.
    \item \textbf{(Modularität)} Für $n=2$, falls $\rho$ zusätzlich ungerade ist
      (d.\,h.\ $\det\rho(\text{komplexe Konjugation}) = -1$), gilt
      $\rho \cong \rho_f$ für eine klassische Modulform $f$.
    \item \textbf{(Automorphie)} Allgemein entspricht $\rho$ einer automorphen
      Darstellung von $\GL_n(\mathbb{A}_\Q)$ vermöge der Langlands-Korrespondenz.
  \end{enumerate}
\end{vermutung}

\begin{bemerkung}
  Die Ungerade-Bedingung für $n=2$ ist notwendig: Das zyklotomische Zeichen
  $\chi_{\mathrm{zyk}}$ ist geometrisch, aber gerade, und entspricht einem
  Hecke-Grössencharakter, nicht einer holomorphen Modulform.
\end{bemerkung}

\begin{vermutung}[Fontaine-Mazur, $n=2$]\label{verm:FM2}
  Jede ungerade, irreduzible, geometrische $2$-dimensionale $p$-adische Galois-Darstellung
  $\rho\colon \GQ \to \GL_2(\Qp)$
  ist modular: Es gibt eine normalisierte spitze Neuform $f$ vom Gewicht $k \geq 1$ mit
  $\rho \cong \rho_f \otimes \Qp$.
\end{vermutung}

% ==============================================================================
\section{Verbindung zum Langlands-Programm}\label{sec:langlands}
% ==============================================================================

Die Fontaine-Mazur-Vermutung ist eng mit dem Langlands-Programm verbunden, das eine
Korrespondenz zwischen
\begin{itemize}
  \item $n$-dimensionalen Darstellungen von $\GQ$ (Galois-Seite) und
  \item automorphen Darstellungen $\pi$ von $\GL_n(\mathbb{A}_\Q)$ (automorphe Seite)
\end{itemize}
vorhersagt.

\begin{satz}[Lokale Langlands-Korrespondenz für $\GL_n$, Harris-Taylor \cite{HT2001},
  Henniart \cite{Henniart2000}]
  Für jede Primzahl $\ell$ gibt es eine lokale Langlands-Korrespondenz:
  \[
    \rec \colon \{ \text{irred.\ glatte Darstell.\ von } \GL_n(\Q_\ell) \}
    \;\longleftrightarrow\;
    \{ \text{irred.\ $n$-dim.\ Weil-Deligne-Darstell.\ von } W_{\Q_\ell} \}.
  \]
\end{satz}

\begin{satz}[$p$-adische lokale Langlands-Korrespondenz für $\GL_2(\Qp)$,
  Colmez \cite{Colmez2010}]
  Es gibt eine $p$-adische lokale Langlands-Korrespondenz für $\GL_2(\Qp)$:
  eine Bijektion (funktoriell, verträglich mit Drehungen und Dualität) zwischen
  \begin{itemize}
    \item absolut irreduziblen $2$-dimensionalen $p$-adischen Darstellungen von
      $G_{\Qp}$, und
    \item absolut irreduziblen unitären zulässigen Banach-Darstellungen von
      $\GL_2(\Qp)$ mit zentralem Zeichen.
  \end{itemize}
\end{satz}

\begin{bemerkung}
  Für $n \geq 3$ ist die $p$-adische lokale Langlands-Korrespondenz noch unbekannt.
  Dies ist eine der Haupthürden für die Fontaine-Mazur-Vermutung bei $n \geq 3$.
\end{bemerkung}

% ==============================================================================
\section{Wiles' Satz und Modularität elliptischer Kurven}\label{sec:wiles}
% ==============================================================================

\begin{satz}[Wiles 1995, vervollständigt durch Breuil-Conrad-Diamond-Taylor
  \cite{BCDT2001}]
  Jede elliptische Kurve $E/\Q$ ist modular: Es gibt eine normalisierte spitze Neuform
  $f \in S_2(\Gamma_0(N))$ mit $L(E,s) = L(f,s)$, äquivalent: $V_p E \cong \rho_f$.
\end{satz}

Der Beweis verläuft in zwei Hauptschritten:
\begin{enumerate}[label=(\roman*)]
  \item \textbf{Residuelle Modularität}: Die Restdarstellung
    $\bar{\rho}_E \colon \GQ \to \GL_2(\F_p)$ (für $p=3$ oder $p=5$) ist entweder
    reduzibel (und damit modular) oder bekannt als modular (via Langlands-Tunnell für
    $\GL_2(\F_3)$).
  \item \textbf{Modularitätslift ($R = \mathbb{T}$)}: Ein deformationstheoretisches
    Argument zeigt, dass der universelle Deformationsring $R$ isomorph zur Hecke-Algebra
    $\mathbb{T}$ ist (Taylor-Wiles-Methode).
\end{enumerate}

\begin{satz}[Taylor-Wiles-Methode \cite{TW1995}]
  Sei $p$ eine ungerade Primzahl und $\bar{\rho}\colon \GQ \to \GL_2(\F_p)$ absolut
  irreduzibel, ungerade und modular. Sei $\rho\colon \GQ \to \GL_2(\Zp)$ eine
  geometrische Deformation von $\bar{\rho}$ mit vorgeschriebenen Hodge-Tate-Gewichten.
  Dann ist $\rho$ modular.
\end{satz}

% ==============================================================================
\section{Kisins Arbeit: Potenziell semistabile Darstellungen}\label{sec:kisin}
% ==============================================================================

\begin{definition_de}[Potenziell semistabil]
  Eine $p$-adische Darstellung $V$ von $G_K$ heißt \emph{potenziell semistabil}, falls
  $V|_{G_{K'}}$ für eine endliche Erweiterung $K'/K$ semistabil wird.
  Nach Bergers Satz ist das äquivalent zu: $V$ ist de Rham.
\end{definition_de}

\begin{satz}[Kisin \cite{Kisin2009}]\label{satz:kisin}
  Sei $p > 2$ und $\bar{\rho}\colon \GQ \to \GL_2(\F_p)$ ungerade, absolut irreduzibel.
  Unter geeigneten technischen Bedingungen (u.\,a.\ $\bar{\rho}|_{G_{\Q(\zeta_p)}}$
  absolut irreduzibel und $\bar{\rho}$ modular) gilt:
  Jede ungerade, irreduzible Deformation $\rho\colon \GQ \to \GL_2(\Qp)$ mit
  $\bar{\rho}_\rho = \bar{\rho}$, die de Rham bei $p$ mit
  Hodge-Tate-Gewichten $0 \leq h_1 < h_2$ ist, ist modular.
\end{satz}

Kisins Schlüsselinnovation sind die \emph{Kisin-Moduln} (auch Breuil-Kisin-Moduln):

\begin{definition_de}[Kisin-Modul]
  Sei $K/\Qp$ eine endliche total verzweigte Erweiterung mit Uniformisierer $\pi$ und
  Eisenstein-Polynom $E$. Ein \emph{Kisin-Modul der Höhe $\leq h$} ist ein endlich
  freier $\mathfrak{S} = W(k)[[u]]$-Modul $\mathfrak{M}$ zusammen mit einer
  $\phi$-semilinearen Abbildung $\phi_\mathfrak{M}\colon \mathfrak{M} \to \mathfrak{M}$,
  deren Kokern von $E(u)^h$ annulliert wird.
\end{definition_de}

\begin{satz}[Kisin \cite{Kisin2006}]
  Der Funktor $\mathfrak{M} \mapsto T(\mathfrak{M})$ von Kisin-Moduln der Höhe $\leq h$
  zu $G_\infty$-stabilen $\Zp$-Gittern ist volltreu und wesentlich surjektiv auf
  Gitter in semistabilen Darstellungen mit Hodge-Tate-Gewichten in $\{0,\ldots,h\}$.
\end{satz}

% ==============================================================================
\section{Emertons komplettierende Kohomologie}\label{sec:emerton}
% ==============================================================================

\begin{definition_de}[Komplettierende Kohomologie, Emerton \cite{Emerton2006}]
  Sei $Y(\Gamma)$ die Modulkurve für eine Kongruenzuntergruppe
  $\Gamma \subset \GL_2(\Z)$.
  Die \emph{komplettierende Kohomologie} ist
  \[
    \widetilde{H}^i := \varprojlim_n \varinjlim_\Gamma H^i(Y(\Gamma), \Z/p^n\Z),
  \]
  die eine $\GL_2(\mathbb{A}_f^p) \times G_{\Qp}$-Wirkung trägt.
\end{definition_de}

\begin{satz}[Emerton \cite{Emerton2010}]\label{satz:emerton}
  Sei $p \geq 5$ und $\bar{\rho}\colon \GQ \to \GL_2(\F_p)$ ungerade, absolut irreduzibel
  und modular. Unter einer ``Generizitäts''-Bedingung an $\bar{\rho}|_{G_{\Qp}}$ gilt:
  Jede ungerade, irreduzible, geometrische Darstellung $\rho\colon \GQ \to \GL_2(\Qp)$
  mit Restdarstellung $\bar{\rho}$ ist modular.
\end{satz}

\begin{bemerkung}
  Gemeinsam decken Kisins Satz~\ref{satz:kisin} und Emertons Satz~\ref{satz:emerton}
  die große Mehrheit der ungeraden, irreduziblen, geometrischen $2$-dimensionalen
  $p$-adischen Galois-Darstellungen ab.
\end{bemerkung}

% ==============================================================================
\section{Taylors potenzielle Modularität}\label{sec:taylor}
% ==============================================================================

\begin{satz}[Taylor \cite{Taylor2002}]\label{satz:taylor}
  Sei $\ell$ eine Primzahl und $\rho\colon \GQ \to \GL_2(\overline{\Q}_\ell)$ ungerade,
  irreduzibel, geometrisch mit Hodge-Tate-Gewichten $\{0, k-1\}$, $k \geq 2$.
  Dann gibt es einen total-reellen Zahlkörper $F$, so dass $\rho|_{G_F}$ modular ist
  (d.\,h.\ einer Hilbertschen Modulform über $F$ entspricht).
\end{satz}

\begin{bemerkung}
  Potenzielle Modularität reicht für viele arithmetische Anwendungen aus:
  \begin{itemize}
    \item meromorphe Fortsetzung und Funktionalgleichung von $L(\rho, s)$;
    \item Ramanujan-Vermutung für $\rho$ (Frobenius-Eigenwerte haben Betrag $p^{(k-1)/2}$);
    \item Sato-Tate-Vermutung (jetzt Satz) für elliptische Kurven über total-reellen
      Körpern \cite{BGHT2011}.
  \end{itemize}
\end{bemerkung}

% ==============================================================================
\section{Offene Probleme}\label{sec:offene_probleme}
% ==============================================================================

Für $n \geq 3$ ist die Fontaine-Mazur-Vermutung nahezu vollständig offen.

\begin{vermutung}[Fontaine-Mazur für $n \geq 3$]\label{verm:FMn3}
  Jede ungerade, irreduzible, geometrische $n$-dimensionale $p$-adische Galois-Darstellung
  $\rho\colon \GQ \to \GL_n(\Qp)$ ist automorph.
\end{vermutung}

Die wichtigsten Hindernisse sind:
\begin{enumerate}[label=(\roman*)]
  \item \textbf{Keine $p$-adische lokale Langlands für $\GL_n$, $n \geq 3$}:
    Colmez' Konstruktion ist wesentlich $\GL_2$-spezifisch.
  \item \textbf{$R = \mathbb{T}$ für $\GL_n$}: Die Taylor-Wiles-Methode erfordert
    ein numerisches Gleichgewicht lokaler Bedingungen, das für $\GL_2$ klappt,
    aber bei $\GL_n$ auf Obstruktionen stößt.
  \item \textbf{Automorphe Formen auf $\GL_n$}: Für $n \geq 3$ gibt es keine direkten
    Analoga zu Shimura-Kurven über $\Q$.
\end{enumerate}

\begin{bemerkung}
  Neuere Entwicklungen -- Scholzes perfektoide Geometrie \cite{Scholze2015},
  prismatische Kohomologie (Bhatt-Scholze \cite{BS2022}) und das kategorielle
  $p$-adische Langlands-Programm (Fargues-Scholze \cite{FS2021}) -- liefern neue Werkzeuge
  für den Angriff auf diese Vermutung.
\end{bemerkung}

\begin{vermutung}[Kategorielle $p$-adische Langlands, Fargues-Scholze \cite{FS2021}]
  Es gibt eine Äquivalenz von $\infty$-Kategorien (nach geeigneter Lokalisierung)
  zwischen kohärenten Garben auf dem Stack der Langlands-Parameter für $\GL_n(\Qp)$
  über der Fargues-Fontaine-Kurve und glatten Darstellungen von $\GL_n(\Qp)$.
\end{vermutung}

% ==============================================================================
\section{Zusammenfassung bekannter Ergebnisse}
% ==============================================================================

\begin{satz}[Stand der Fontaine-Mazur-Vermutung für $n=2$]
  Die folgenden Fälle von Vermutung~\ref{verm:FM2} sind bewiesen:
  \begin{enumerate}[label=(\roman*)]
    \item Alle elliptischen Kurven $E/\Q$: $V_pE$ ist modular \cite{BCDT2001}.
    \item Ungerade, irreduzible $p$-adische Darstellungen mit Hodge-Tate-Gewichten
      in $\{0,1\}$ und residuell modular (unter milden lokalen Bedingungen bei $p$)
      \cite{Kisin2009, Emerton2010}.
    \item Alle solchen Darstellungen werden über einem total-reellen Körper modular
      \cite{Taylor2002}.
    \item Gewicht-$1$-Modulformen entsprechen ungeraden Artin-Darstellungen
      (Deligne-Serre \cite{DS1974}).
  \end{enumerate}
\end{satz}

% ==============================================================================
\section{Fazit}\label{sec:fazit}
% ==============================================================================

Die Fontaine-Mazur-Vermutung zählt zu den zentralen offenen Problemen der arithmetischen
Geometrie. Sie verbindet drei tiefe mathematische Theorien:
$p$-adische Hodge-Theorie (Fontaines Periodenringe), Galois-Darstellungen und das
Langlands-Programm. Der Fall $n=2$ ist durch die kombinierte Arbeit von Wiles, Kisin
und Emerton weitgehend geklärt: Nahezu alle ungeraden irreduziblen geometrischen
$2$-dimensionalen $p$-adischen Darstellungen sind modular.
Für $n \geq 3$ stellt die Vermutung eine der großen Herausforderungen der aktuellen
Forschung dar, für die neuartige Werkzeuge -- etwa Scholzes perfektoide Geometrie
und das Fargues-Fontaine-Programm -- neue Hoffnung geben.

Die vollständige Lösung der Fontaine-Mazur-Vermutung würde einen Meilenstein in der
Zahlentheorie darstellen und insbesondere die Automorphizität aller $\ell$-adischen
Kohomologiegruppen glatter projektiver Varietäten über $\Q$ implizieren.

\begin{thebibliography}{99}

\bibitem{BCDT2001}
C.\ Breuil, B.\ Conrad, F.\ Diamond, R.\ Taylor,
\textit{On the modularity of elliptic curves over $\mathbf{Q}$: wild $3$-adic exercises},
J.\ Amer.\ Math.\ Soc.\ \textbf{14} (2001), Nr.\ 4, 843--939.

\bibitem{Berger2002}
L.\ Berger,
\textit{Représentations $p$-adiques et équations différentielles},
Invent.\ Math.\ \textbf{148} (2002), Nr.\ 2, 219--284.

\bibitem{BGHT2011}
T.\ Barnet-Lamb, D.\ Geraghty, M.\ Harris, R.\ Taylor,
\textit{A family of Calabi-Yau varieties and potential automorphy II},
P.\ R.\ I.\ M.\ E.\ S.\ \textbf{47} (2011), Nr.\ 1, 29--98.

\bibitem{BS2022}
B.\ Bhatt, P.\ Scholze,
\textit{Prisms and prismatic cohomology},
Ann.\ of Math.\ (2) \textbf{196} (2022), Nr.\ 3, 1135--1275.

\bibitem{Colmez2010}
P.\ Colmez,
\textit{Représentations de $\mathrm{GL}_2(\mathbf{Q}_p)$ et $(\phi, \Gamma)$-modules},
Astérisque \textbf{330} (2010), 281--509.

\bibitem{Deligne1971}
P.\ Deligne,
\textit{Formes modulaires et représentations $\ell$-adiques},
Séminaire Bourbaki, vol.\ 1968/69, Exposé 355, Lect.\ Notes in Math.\ Bd.\ 179,
Springer (1971), 139--172.

\bibitem{DS1974}
P.\ Deligne, J.-P.\ Serre,
\textit{Formes modulaires de poids $1$},
Ann.\ Sci.\ École Norm.\ Sup.\ (4) \textbf{7} (1974), 507--530.

\bibitem{Emerton2006}
M.\ Emerton,
\textit{On the interpolation of systems of eigenvalues attached to automorphic Hecke eigenforms},
Invent.\ Math.\ \textbf{164} (2006), Nr.\ 1, 1--84.

\bibitem{Emerton2010}
M.\ Emerton,
\textit{Local-global compatibility in the $p$-adic Langlands programme for $\mathrm{GL}_{2/\mathbf{Q}}$},
Preprint (2010).

\bibitem{FM1995}
J.-M.\ Fontaine, B.\ Mazur,
\textit{Geometric Galois representations},
in \textit{Elliptic Curves, Modular Forms, \& Fermat's Last Theorem (Hong Kong, 1993)},
Int.\ Press, Cambridge, MA (1995), 41--78.

\bibitem{Fontaine1979}
J.-M.\ Fontaine,
\textit{Modules galoisiens, modules filtrés et anneaux de Barsotti-Tate},
Astérisque \textbf{65} (1979), 3--80.

\bibitem{Fontaine1982}
J.-M.\ Fontaine,
\textit{Sur certains types de représentations $p$-adiques du groupe de Galois d'un corps local},
Ann.\ of Math.\ (2) \textbf{115} (1982), Nr.\ 3, 529--577.

\bibitem{Fontaine1994}
J.-M.\ Fontaine,
\textit{Le corps des périodes $p$-adiques},
Astérisque \textbf{223} (1994), 59--111.

\bibitem{FS2021}
L.\ Fargues, P.\ Scholze,
\textit{Geometrization of the local Langlands correspondence},
Preprint (2021), arXiv:2102.13459.

\bibitem{HT2001}
M.\ Harris, R.\ Taylor,
\textit{The Geometry and Cohomology of Some Simple Shimura Varieties},
Princeton Univ.\ Press (2001).

\bibitem{Henniart2000}
G.\ Henniart,
\textit{Une preuve simple des conjectures de Langlands pour $\mathrm{GL}(n)$ sur un corps $p$-adique},
Invent.\ Math.\ \textbf{139} (2000), Nr.\ 2, 439--455.

\bibitem{Kisin2006}
M.\ Kisin,
\textit{Crystalline representations and $F$-crystals},
Progr.\ Math.\ \textbf{253}, Birkhäuser (2006), 459--496.

\bibitem{Kisin2009}
M.\ Kisin,
\textit{Moduli of finite flat group schemes, and modularity},
Ann.\ of Math.\ (2) \textbf{170} (2009), Nr.\ 3, 1085--1180.

\bibitem{Scholze2015}
P.\ Scholze,
\textit{On torsion in the cohomology of locally symmetric varieties},
Ann.\ of Math.\ (2) \textbf{182} (2015), Nr.\ 3, 945--1066.

\bibitem{Taylor2002}
R.\ Taylor,
\textit{Remarks on a conjecture of Fontaine and Mazur},
J.\ Inst.\ Math.\ Jussieu \textbf{1} (2002), Nr.\ 1, 125--143.

\bibitem{Taylor2004}
R.\ Taylor,
\textit{Galois representations},
Ann.\ Fac.\ Sci.\ Toulouse Math.\ (6) \textbf{13} (2004), Nr.\ 1, 73--119.

\bibitem{TW1995}
R.\ Taylor, A.\ Wiles,
\textit{Ring-theoretic properties of certain Hecke algebras},
Ann.\ of Math.\ (2) \textbf{141} (1995), Nr.\ 3, 553--572.

\bibitem{Wiles1995}
A.\ Wiles,
\textit{Modular elliptic curves and Fermat's last theorem},
Ann.\ of Math.\ (2) \textbf{141} (1995), Nr.\ 3, 443--551.

\end{thebibliography}

\end{document}
